\begin{tikzpicture}[
  >={Latex[length=3mm,width=2.2mm]},
  every node/.style={font=\small},
  ivec/.style={->,line width=1pt,vpurple},
  bvec/.style={->,line width=1.2pt,bblue},
  Fvec/.style={->,line width=1.2pt,black}
]

% --- In alto: correnti concordi (si attraggono) ---
\begin{scope}
  % Filo 1 (grigio inclinato)
  \fill[blobgray!60,draw=black!50,rotate=20]
    (-2.5,-0.1) rectangle (0.5,0.1);
  % Filo 2 (grigio inclinato)
  \fill[blobgray!60,draw=black!50,rotate=20]
    (1.5,-0.1) rectangle (4.5,0.1);
  % Correnti
  \draw[ivec] (-1.8,0.0) -- ++(-0.5,-0.2) node[above left] {$i_1$};
  \draw[ivec] (2.8,1.1) -- ++(0.5,0.2) node[above right] {$i_2$};
  % B
  \draw[bvec] (-1.0,0.6) -- ++(0,1.0) node[above] {$\mathbf{B}_2$};
  \draw[bvec] (2.5,0.5) -- ++(0,-1.0) node[below] {$\mathbf{B}_1$};
  % F (attrazione)
  \draw[Fvec] (-0.5,0.3) -- ++(1.0,0) node[right] {$\mathbf{F}$};
  \draw[Fvec] (2.0,0.8) -- ++(-1.0,0) node[left] {$-\mathbf{F}$};
  % r
  \draw[<->,thin] (0.2,-0.2) -- (1.8,0.4);
  \node[below] at (1.0,0.1) {$r$};
\end{scope}

% --- In basso: correnti discordi (si respingono) ---
\begin{scope}[yshift=-3.5cm]
  \fill[blobgray!60,draw=black!50,rotate=0]
    (-2.5,-0.1) rectangle (0.0,0.1);
  \fill[blobgray!60,draw=black!50,rotate=0]
    (2.0,-0.1) rectangle (4.5,0.1);
  % Correnti (opposte)
  \draw[ivec] (-1.8,0.2) -- ++(-0.5,0) node[above left] {$i_1$};
  \draw[ivec] (3.5,0.2) -- ++(0.5,0) node[above right] {$i_2$};
  % B
  \draw[bvec] (-0.5,-0.2) -- ++(0,-1.0) node[below] {$\mathbf{B}_2$};
  \draw[bvec] (2.5,-0.2) -- ++(0,-1.0) node[below] {$\mathbf{B}_1$};
  % F (repulsione)
  \draw[Fvec] (-0.5,0.0) -- ++(-1.0,0) node[left] {$\mathbf{F}$};
  \draw[Fvec] (2.5,0.0) -- ++(1.0,0) node[right] {$-\mathbf{F}$};
  % r
  \draw[<->,thin] (0.0,-0.5) -- (2.0,-0.5);
  \node[below] at (1.0,-0.5) {$r$};
\end{scope}

\end{tikzpicture}
